%======================================================================
\section{Substructural Logics}\label{sec:fl}

In this section we extend the translation to richer languages and show
that the analogue of Corollary~\ref{cor:many-logics} holds over
$\FLe$ in the multiplicative language and over $\FLew$
in the full language, but not over $\FLe$ with additive
connectives. The results of this section are not used elsewhere in the
paper.

\subsection{The translation in a richer language}

Let $\mathcal L$ be a language that contains $\to$ and possibly some of the
binary connectives $\cdot$ (fusion), $\wedge$ and $\vee$ and some of the
constants $\bot,\top,0,1$. \AP An \intro{atom} is a variable or a constant. \AP We
consider the full Lambek calculus with exchange $\intro*\FLe$ and its
extension $\intro*\FLew$ by weakening, in the language
$\{\to,\cdot,\wedge,\vee,0,1\}$, possibly expanded by $\bot$ and $\top$
(see~\cite{GJKO07}). Intuitionistic logic in this language contains
$\FLew$. The full Lambek calculus $\intro*\FL$ without exchange has two
implications $\backslash$ and $/$ in place of $\to$. Logics in $\mathcal L$, the sets $\Gen{\alpha}$ and
occurrences are defined as in Sections~\ref{sec:prelim}
and~\ref{sec:proof}, except that an occurrence $u$ is a \kl{leaf} if
$\subf{\alpha}{u}$ is an \kl{atom}, and if $\subf{\alpha}{u}=\gamma\circ\delta$ for a binary
connective $\circ$, then $\gamma$ is at $u0$ and $\delta$ is at $u1$.

The translation $\St$ is defined as in Definition~\ref{def:S}, with the
following changes. If $\subf{\alpha}{u}=\subf{\alpha}{u0}\circ\subf{\alpha}{u1}$ with
$\circ\in\{\cdot,\wedge,\vee\}$, then both $u0$ and $u1$ have the same
\kl{polarity} as $u$. For a \kl{leaf} $u$ we put $\rv{u}=\subf{\alpha}{u}$, so that constants
are treated like variables, and for a \kl{non-leaf occurrence} $u$ with
$\subf{\alpha}{u}=\subf{\alpha}{u0}\circ\subf{\alpha}{u1}$ we put $\body{u}=\rv{u0}\circ \rv{u1}$.

\begin{theorem}\label{thm:main-fl}
Let $L$ be a \kl{logic} in $\mathcal L$ containing $\axI$. If $\alpha$ is \kl{minimal}
in $L$ and $\St(\alpha)\in L$, then $\St(\alpha)$ is \kl{minimal} in $L$.
\end{theorem}
\begin{proof}
The proof of Theorem~\ref{thm:main} in Section~\ref{sec:proof} goes through after the following changes:
\begin{itemize}
\item variables of $\alpha$ are replaced by \kl{atoms} of $\alpha$, except
  in the condition on $\delta$ in the definition of $\Gen{\alpha}$, which still
  concerns variables only;
\item in the ``only if'' part of the proof of
  Proposition~\ref{cor:criterion} (Appendix~\ref{app:criterion}), occurrences of $\to$ are replaced by
  occurrences of connectives and constants;
\item in Lemmas~\ref{lem:shape}(d) and~\ref{lem:key},
  the implication in $\body{u}$ is replaced by the main connective of the
  corresponding occurrence.\qedhere
\end{itemize}
\end{proof}

\subsection{Soundness over \texorpdfstring{$\FLe$}{FLe} and \texorpdfstring{$\FLew$}{FLew}}

\begin{lemma}\label{lem:sound-fl}
In each of the following cases, $\alpha\to\St(\alpha)\in L$ for every
formula $\alpha$ of $\mathcal L$:
\begin{enumerate}
\item[(i)] $\mathcal L\subseteq\{\to,\cdot,0,1\}$ and $L\supseteq\FLe$;
\item[(ii)] $L\supseteq\FLew$; in particular, $L\supseteq\IL$.
\end{enumerate}
\end{lemma}
\begin{proof}
We argue as in the proof of Lemma~\ref{lem:sound}, in the sequent calculus
for $\FLe$, with weakening in case~(ii). The case of fusion is
similar to that of implication, and again every formula of $\Gamma_u$ is
used exactly once; this proves (i). For $\wedge$ and $\vee$, the rules
$({\wedge}\mathrm{R})$ and $({\vee}\mathrm{L})$ share their context between
the two upper sequents, and in each upper sequent we remove the formulas $\dprem{v}$
with $v$ in the other subtree by weakening; this proves (ii).
\end{proof}

\begin{corollary}\label{cor:fl}
In each case of Lemma~\ref{lem:sound-fl}, $L\lplus \alpha=L\lplus \St(\alpha)$ for every
formula $\alpha$ of $\mathcal L$. If, moreover, $\alpha$ is \kl{minimal} in $L$,
then so is $\St(\alpha)$.
\end{corollary}
\begin{proof}
As in the proof of Corollary~\ref{cor:many-logics},
using Lemma~\ref{lem:sound-fl} and Theorem~\ref{thm:main-fl}.
\end{proof}

\subsection{Additive connectives over \texorpdfstring{$\FLe$}{FLe}}\label{sec:additive}

The proof of Theorem~\ref{thm:main-fl} uses only (C1)--(C3), and \kl{polarity}
enters only through Lemma~\ref{lem:sound-fl}. It is exactly this lemma that
fails over $\FLe$ in the presence of additive connectives. The
following two observations, both verified with a cut-free prover for
$\FLe$ (\texttt{tools/fle\_check.py}), show that the difficulty
cannot be removed simply by replacing each $\dprem{u}$ by
$\dprem{u}\wedge1$ as Jeřábek does.

\begin{remark}\label{rem:wrapper}
Let $\mathcal L=\{\to,\cdot,\wedge,\vee,0,1\}$.
\begin{enumerate}
\item[(a)] \emph{Without the conjunctions $\wedge1$, $\St$ is not sound over
  $\FLe$.} The formula $\alpha=(q\to q)\wedge(r\to r)$ is \kl{minimal}
  in $\FLe$, but
  \[
  \St(\alpha)=((q\to q)\to \evfp{q\to q})\to((r\to r)\to \evfp{r\to r})
  \to((\evfp{q\to q}\wedge \evfp{r\to r})\to\evfp{\alpha})\to \evfp{\alpha}
  \]
  is not provable in $\FLe$: each upper sequent of the rule
  $({\wedge}\mathrm{R})$ would have to discard one of the \kl{antecedents}
  $(q\to q)\to\evfp{q\to q}$ and $(r\to r)\to\evfp{r\to r}$.
\item[(b)] \emph{With the conjunctions $\wedge1$, minimality is always lost.} Jeřábek's
  $\varphi^+$ has $\dprem{u}\wedge1$ in place of $\dprem{u}$, which makes it sound over
  $\FLe$ \cite[Thm.~3.8]{Jer16}. The factor $\dprem{\rt}\wedge1$
  is never discarded in the proof, so replacing its conjunct $1$ by $z$
  gives an element of $\Gen{\varphi^+}$ that is still provable, by
  $({\wedge}\mathrm{L})$. Hence $\varphi^+$ is not \kl{minimal} in any \kl{logic}
  $L\supseteq\FLe$, including $\FLew$, $\IL$ and $\CL$,
  whenever $\varphi$ is compound. For example, for $\varphi=p\to p$ we have
  $\varphi^+=(((p\to p)\to \evfp{p\to p})\wedge 1)\to \evfp{p\to p}$, and
  $(((p\to p)\to \evfp{p\to p})\wedge z)\to \evfp{p\to p}$ is provable.
\end{enumerate}
\end{remark}

Exhaustive tests over all formulas in $\to,\wedge,\vee$ with at most four
binary connectives and three variables give the following picture. Up to
renaming, there are $413$ formulas \kl{minimal} in $\FLe$. For $164$ of
them $\St(\alpha)$ is provable, and in all these cases it is \kl{minimal}, as
predicted by Theorem~\ref{thm:main-fl}. The formula $\varphi^+$ is provable in
all $413$ cases and \kl{minimal} in none. Replacing $\dprem{u}$ by $\dprem{u}\wedge1$ only for
occurrences $u$ \kl{strictly below} an additive occurrence yields a provable formula
in all $413$ cases, but a \kl{minimal} one in only $255$. For
$\alpha=p\vee(q\to q)$, for instance, the conjunction $\wedge1$ is not needed because
the rule $({\vee}\mathrm{R})$ selects one disjunct. Thus a
minimality-preserving translation over $\FLe$, if it exists, must
choose the conjunctions $\wedge1$ according to the actual proof. We leave open whether such a
translation exists, as well as the non-commutative case $\FL$.
